\chapter{Benign Tumours}

\section*{Definition and Overview}
A benign tumour is an abnormal growth of cells that is not cancerous. Unlike malignant (cancerous) tumours, benign tumours grow slowly, remain confined to their site of origin, and do not invade surrounding tissues or spread (metastasise) to distant parts of the body. They can arise in almost any tissue and usually cause problems only by growing large or by pressing on nearby structures.

Most benign tumours are harmless, and many are found incidentally during investigations for unrelated complaints. A small minority, however, can cause significant symptoms, produce hormones, bleed, or carry a low risk of transforming into cancer over many years. Distinguishing a benign tumour from cancer reliably requires medical assessment and, in many cases, imaging or tissue sampling.

\section*{Epidemiology}
Benign tumours are extremely common, and their frequency generally increases with age. Because most never cause trouble, the true population figures are not known, and many are detected only when a scan or operation is performed for another reason.

\begin{itemize}
    \item \textbf{Common types:} lipomas (fatty lumps under the skin), skin moles (naevi), breast fibroadenomas, uterine fibroids, and adenomatous polyps of the bowel
    \item \textbf{Sex and age patterns:} uterine fibroids are common in women of reproductive age, while bowel polyps become increasingly frequent after the age of 50
    \item \textbf{Intracranial tumours:} benign brain tumours such as meningiomas and pituitary adenomas are uncommon but occur mainly in adults
    \item \textbf{Screening detection:} bowel cancer screening programmes regularly detect benign polyps that would otherwise have gone unnoticed
\end{itemize}

\section*{Aetiology and Pathophysiology}
Benign tumours arise when a single cell acquires a genetic change that allows it to divide more frequently than normal. The resulting mass consists of cells that closely resemble the tissue of origin and that remain localised, because they lack the properties that enable cancer cells to invade neighbouring structures or enter the bloodstream and lymphatics.

Most genetic changes that give rise to benign tumours are acquired during life rather than inherited, and the reasons they occur are usually unclear. Some benign tumours are influenced by hormones, such as uterine fibroids, which grow during reproductive life and shrink after the menopause. A minority of benign tumours secrete hormones or other biologically active substances, producing symptoms distant from the site of the growth. A small number, such as certain bowel polyps, have the capacity to transform into malignant lesions over many years, which is the rationale for their removal or surveillance.

\section*{Clinical Features}
Many benign tumours produce no symptoms at all and are discovered by chance. When symptoms occur, they typically reflect the size and location of the growth.

\begin{itemize}
    \item \textbf{A palpable lump:} a slow-growing swelling or nodule that may be felt under the skin or within a breast or gland
    \item \textbf{Pressure effects:} symptoms from compression of nearby organs, such as difficulty breathing, swallowing, or passing urine
    \item \textbf{Bleeding:} for example a bowel polyp causing blood in the stools, or a uterine fibroid causing heavy or painful periods
    \item \textbf{Hormone effects:} hormone-secreting tumours can cause high blood pressure, excessive sweating, or changes in body appearance, as with some pituitary or adrenal tumours
    \item \textbf{Neurological symptoms:} brain and spinal tumours may cause headache, seizures, or focal weakness and numbness
    \item \textbf{Typical absence:} benign tumours usually do not cause unexplained weight loss, fever, or night sweats, which are more suggestive of malignancy or infection
\end{itemize}

\section*{Differential Diagnosis}
A lump or abnormal growth can have several causes, and the principal task is to distinguish a benign tumour from cancer.

\begin{itemize}
    \item \textbf{Malignant tumours (cancer):} the most important consideration for any new, growing, or changing lump
    \item \textbf{Cysts:} fluid-filled sacs, such as sebaceous cysts, ovarian cysts, and Baker's cysts
    \item \textbf{Abscesses:} localised collections of pus resulting from infection
    \item \textbf{Enlarged lymph nodes:} commonly due to infection or inflammation, but sometimes due to lymphoma or metastatic cancer
    \item \textbf{Inflammatory masses:} such as those occurring in inflammatory bowel disease, sarcoidosis, or granulomatous conditions
    \item \textbf{Pre-cancerous lesions:} such as some bowel polyps and dysplastic moles, which may progress to malignancy
    \item \textbf{Hernias and other anatomical swellings:} which can mimic a tumour on examination
\end{itemize}

\section*{When to Seek Medical Care}
Any new lump, or a change in a known lump, should be assessed by a doctor. Seek prompt review if a lump:

\begin{itemize}
    \item Is growing quickly or changing in size, shape, texture, or feel
    \item Becomes painful, firm, irregular, or fixed to the surrounding tissue
    \item Causes new bleeding, discharge, or symptoms of pressure on nearby organs
    \item Is accompanied by unexplained weight loss, fever, or night sweats
    \item Changes in colour or outline, particularly a skin mole, which should be checked promptly
\end{itemize}

\textbf{Call 000 (or local emergency services) for:}
\begin{itemize}
    \item Sudden severe headache or collapse in a person with a known brain tumour
    \item Sudden weakness, seizures, or confusion
    \item Severe breathing difficulty or the inability to swallow
    \item Heavy or uncontrolled bleeding
\end{itemize}

\section*{Diagnosis and Investigations}
\begin{itemize}
    \item \textbf{History and examination:} assessment of the lump's size, texture, mobility, tenderness, and relationship to surrounding structures, together with symptoms and risk factors
    \item \textbf{Imaging:} ultrasound, X-ray, CT, or MRI to define the size, location, and internal characteristics of the growth
    \item \textbf{Biopsy:} removal of a sample of tissue for examination under a microscope, which is the most reliable way to determine whether a tumour is benign or malignant
    \item \textbf{Endoscopy:} direct visualisation of growths within hollow organs, such as colonoscopy for bowel polyps or cystoscopy for bladder lesions
    \item \textbf{Blood tests:} including hormone assays for suspected endocrine tumours and baseline blood counts where bleeding has occurred
\end{itemize}

\section*{Management}
Management depends on the type of tumour, its location and size, and whether it is causing symptoms or carries a risk of transformation.

\begin{itemize}
    \item \textbf{Observation:} many benign tumours require no treatment and are simply monitored with periodic review and interval imaging
    \item \textbf{Surgical removal:} recommended when a tumour causes symptoms, is enlarging, or has the potential to become malignant
    \item \textbf{Endoscopic removal:} for polyps or small growths within the bowel, stomach, bladder, or airway
    \item \textbf{Medication:} used for certain types, such as hormonal therapy for uterine fibroids, or dopamine agonists for some pituitary tumours
    \item \textbf{Radiotherapy or radiosurgery:} selected treatment for some benign brain tumours that are not surgically removable
    \item \textbf{Follow-up:} regular surveillance to detect growth or transformation early where that risk exists
\end{itemize}

\section*{Complications}
\begin{itemize}
    \item \textbf{Compression and obstruction:} pressure on or blockage of nearby structures, such as a polyp obstructing the bowel or a fibroid compressing the bladder or ureters
    \item \textbf{Bleeding:} for example from bowel polyps or uterine fibroids, which can cause anaemia
    \item \textbf{Torsion:} twisting of a stalked growth, such as some lipomas or ovarian masses, causing sudden pain
    \item \textbf{Hormonal effects:} endocrine tumours can produce excess hormones with systemic consequences
    \item \textbf{Cosmetic and psychological impact:} visible lumps may cause embarrassment or anxiety
    \item \textbf{Malignant transformation:} a low risk in a minority of tumour types, such as certain bowel polyps
    \item \textbf{Treatment-related complications:} including the risks of surgery, anaesthesia, and invasive procedures
\end{itemize}

\section*{Prognosis and Outcomes}
The outlook for benign tumours is generally excellent. Most do not shorten life, and many require no treatment at all.

\begin{itemize}
    \item Benign tumours do not spread to distant organs and are not directly life-threatening in the majority of cases
    \item Where problems occur, they usually result from the tumour's location, size, hormonal activity, or bleeding
    \item Surgical or endoscopic removal is usually curative for those that require treatment, although some types can recur locally
    \item A small number of tumour types carry a low risk of becoming malignant over time, which is why they are removed or followed up
    \item Regular surveillance provides reassurance and enables early intervention if a change is detected
\end{itemize}

\section*{Prevention and Self-Care}
Most benign tumours cannot be prevented, but awareness and early detection reduce the chance of complications.

\begin{itemize}
    \item Participate in bowel cancer screening, which can find and remove benign polyps before they become malignant
    \item Examine the skin regularly and seek advice about new or changing moles
    \item Practise breast self-awareness and attend recommended breast screening
    \item See a doctor promptly about any new, growing, or changing lump
    \item Attend scheduled follow-up for known benign tumours so that any change is detected early
    \item Maintain a healthy lifestyle, since weight and physical activity influence the risk of some hormone-sensitive growths
\end{itemize}

\section*{Sources and Further Reading}
The following organisations provide reliable information on benign tumours for patients, families, and health professionals.

\begin{itemize}
    \item NHS. \textit{Information on non-cancerous growths, lumps, and their assessment.} https://www.nhs.uk/conditions/
    \item Mayo Clinic. \textit{Overview of benign tumours, including diagnosis and treatment.} https://www.mayoclinic.org/
    \item MSD Manual. \textit{Professional information on benign versus malignant tumours.} https://www.msdmanuals.com/
    \item MedlinePlus. \textit{Health information on tumours and abnormal growths for patients.} https://medlineplus.gov/
    \item National Cancer Institute. \textit{Educational resources on the biology and classification of tumours.} https://www.cancer.gov/
\end{itemize}
